\hypertarget{source-terminal-period}{}
Use the column lattice \(N_u\) and the integrality tests on the row lattice \(M_\mu\) of Section~\ref{subsec:selected-period}. We describe
the effect of the class sums on the column and the terminal
reading on these tests. These steps take place over the
coefficient field \(K\). Write \(\mathcal O_w\) for its local
ring and 
\[
  v=v_w/e_w
\]
 normalized by \(v(p)=1\). Let \(D_{\Gamma,u}\) be the
common stabilizer multiple of the chain and put
\(\Delta_w=v_p(D_{\Gamma,u})\). The multiple includes the
terminal classes of the reference and used twisted chains.

\begin{proposition}[The terminal valuation without a neighbor operator]\label{prop:unserved-terminal-valuation}
At a level-free place without the local orbit condition,
with \(p\notin Q(\mu(w))\), suppose the stage coefficients
have normalized valuations at least \(B_j\) in the actual
incoming frames. For each terminal period value \(b\) on the local tests of \(M_\mu\),
\[
 v_p(b)\ge\sum_{j=1}^m B_j-m\Delta_w,\qquad \Delta_w=v_p(D_{\Gamma,u}),
\]
where \(D_{\Gamma,u}\) is the common
chain stabilizer multiple including both terminal class sets.
\end{proposition}

\begin{proof}
\proofstep{Pass from coefficient valuations to outgoing integral Weyl tests.}
At a level-free place the localized lattice of stage \(i\)
is read by the local Weyl values at each of its classes.
In the coordinates of the incoming member of the chain these
values are integral for every flag-minor test. Conversely,
for a vector in the rational stage space, integrality of
these values is the defining local lattice condition.

For
one incoming and one outgoing class the kernel reading is
the derivative polynomial in the incoming frame, multiplied
by a unit. Expand it in products of flag minors on the
incoming frame. Integral straightening as in
Section~\ref{subsec:cycle-jets} shows that a coefficient
bound \(B_j\) gives coefficients of valuation at least \(B_j\)
also in such a reading.
Each straightened coefficient is an integral linear combination; explicitly,
\[
  v(a_i)\ge B_j,\quad u_i\in{\cal O}_w\quad\Longrightarrow\quad v\!\left(\sum_i u_i a_i\right)\ge B_j .
\]
 Applying the incoming local
functional gives the same bound on the outgoing values.
The class map is rational and the stage space is stable
under the full operator. These two facts let one use the
local lattice criterion after rescaling the output. Indeed
a nonzero rational value of valuation at least \(B_j\)
has 
\[
  v_w(y)=e_w v(y)\ \ge\ \lceil e_w B_j\rceil\qquad(y\ne0,\ v(y)\ge B_j),
\]
 so the rescaling can be
by a power of a uniformizer chosen in \(K\). This applies
also when \(B_j<0\).

\proofstep{Charge an incoming stabilizer at each stage where it occurs.}
The full map is the sum of the class maps with weights
\(1/|\Gamma_c|\) on the incoming class \(c\). This costs
at most \(\Delta_w\) in \(v\) at stage \(j\) when \(j>1\).
The first incoming group is trivial for the reference
and every used twist. Consequently the accumulated
valuation of the stage-\(m\) vector, relative to its
localized stage lattice, is at least
\[
           \sum_{j=1}^m B_j-(m-1)\Delta_w .
\]
The same composition applies to a bound proved on
the entire column at a specified level with its specified
starting lattice. Products of the successive scalars add
their valuations, so composition introduces no rounding
loss. The estimates without the local orbit condition start in the ramified
submodule of the original trace-zero column. The
calculation at the prime assigned to the sequence uses the lattice with its extra
local condition at the assigned prime.

For a vector \(y\) in the localized stage-\(m\) lattice,
first take its image under the point restriction and the
lowering map into the last compact row model. The point
lattice comparison sends integral stage vectors into the
designated integral terminal lattice. For the reference
chain, the normalized Petersson period decomposes over
terminal classes as
\begin{equation}\label{eq:terminal-period-valuation}
     P_\mu(y,F)=
         \sum_{c}\frac{q_{\mu,c}(y,F)}{|\Gamma_{\mu,c}|}.
\end{equation}
The same formula holds for the used twist with its own
terminal class restrictions. It follows by restricting
both functions to each orbit in the finite double quotient.
The characteristic projectors of the classes are disjoint.

In the invariant Hermitian pairing each class contributes
the local pairing with weight \(1/|\Gamma_{\mu,c}|\).
Writing \(P_c\) for the period restricted to class \(c\) and \(\langle\ ,\ \rangle_c\) for its local pairing, this weighting reads
\[
  P_c(y,F)=|\Gamma_{\mu,c}|^{-1}\langle y_c,F_c\rangle_c .
\]
The real type pairing uses the orthonormal frame for the
positive terminal form and the normalized invariant product
in that frame. The same normalization is used for the
Weyl tests in the terminal lattice.

The integral row condition represents the restriction of
each sharp test \(F\) by a flag-minor polynomial on the
local terminal lattice at \(c\). This is the same lattice
whose derivative dual reads the point image of the
column. More explicitly, in a lattice frame take the
span of the flag-minor products of the prescribed diagram.
A change of basis of the underlying lattice and the
inverse-transpose change in its dual transport the
constant-coefficient derivative pairing. They carry
the row span and the dual column lattice together.
In the reference frame the numerator of the pairing for a row monomial \(X^\alpha\)
is given by \(\partial^\alpha\) at zero.

For monomials this numerator has the diagonal form
\[
  (\partial^\alpha X^\beta)(0)=\begin{cases}\alpha!&\alpha=\beta,\\0&\alpha\ne\beta.\end{cases}
\]
 Divide by the product
of the squared Fock norms of the highest polynomials, as in
Section~\ref{subsec:terminal-realization}. The distinguished
real normalized form is this polynomial pairing on the
restricted representation, with the row conjugation in
its coordinates. Thus evaluating an integral point
column on the integral row span gives the embedded
value of its \(K\)-valued Weyl pairing. Equivalently
multiply the period of the two class restrictions by
\(|\Gamma_{\mu,c}|\) to recover that pairing, since the
single-class weight in the real period is its reciprocal.

In particular, if \(P_c\) denotes the period of the two restrictions to this class,
\[
  q_{\mu,c}(y,F)=|\Gamma_{\mu,c}|P_c .
\]
This gives \(q_{\mu,c}(y,F)\) in \(K\) of valuation
at least zero. The same duality applies in the used
twisted chain with its transported row span. Its
integral period reading is used at \(p\notin Q(\mu)\).
The claim holds for localized \(y\) by \(K\)-linearity in
the column. For the normalized integral column, the scalar
\(D_{\Gamma,u}/|\Gamma_{\mu,c}|\) is integral for every class \(c\).
Consequently \eqref{eq:terminal-period-valuation} gives
\[
 D_{\Gamma,u}P_\mu(y,F)
   =\sum_c \frac{D_{\Gamma,u}}{|\Gamma_{\mu,c}|}\,q_{\mu,c}(y,F)
        \in\mathcal O_w,\qquad v(P_\mu(y,F))\ge-\Delta_w .
\]
This incurs one stabilizer loss for the entire terminal sum. A scalar removed from the column before this
reading multiplies the period and adds its own valuation.

It follows for the case without the local orbit condition that the valuation of the
rational period is at least
\begin{equation}\label{eq:unserved-terminal-bound}
                   \sum_{j=1}^m B_j-m\Delta_w .
\end{equation}
For a place with the local orbit condition the same terminal loss removes the
single \(+\Delta_w\) in the column bound. At the prime assigned to the sequence the stabilizer multiple is prime to \(T_u\), so the
terminal reading preserves the bound at the assigned prime.
Indeed the terminal denominator at this prime has
\[
  T_u\nmid D_{\Gamma,u}\quad\Longrightarrow\quad v_{T_u}(D_{\Gamma,u})=0 .
\]
 At
either type prime the local transport estimate is
already a bound on the period against the row test.

To apply this to the actual test space for \(w\), take
\(\mu=\mu(w)\). The level exclusion \(p\notin Q(\mu(w))\)
refers to the coefficient place of this reading. The row space at the refined level image of these tests, localized at \(w\), is \(R_w\).
The selected period comparison over \(K\) induces the same
valuation on the common rational pairing. The local bound
therefore applies to every test vector in \(R_w\), including
localized linear combinations of the images.

Section~\ref{subsec:terminal-realization} identifies these
classwise row spaces at the refined levels and terminal points on both tuples.
\label{cor:terminal-period-1}
\end{proof}

